\documentclass[12pt]{article}
\usepackage{amsmath, amssymb, geometry}
\geometry{a4paper, margin=1in}

\title{Notes: Technological Advancements and Automation}
\author{}
\date{}

\begin{document}

\maketitle

\section*{Exercise 1: Labor Productivity Growth}
\textbf{Scenario:} A factory introduces advanced robotics in its production line, reducing human labor requirements but increasing overall output. Before automation, 100 workers produced 10,000 units per year. After automation, the factory produces 25,000 units annually with only 50 workers.

\textbf{Questions:}
\begin{enumerate}
    \item Calculate labor productivity (output per worker) before and after automation.
    \item What is the percentage increase in labor productivity due to automation?
    \item Discuss how this improvement might affect wages and employment in the factory.
\end{enumerate}

\section*{Exercise 2: The Impact of Automation on Employment}
\textbf{Scenario:} In an industry, technological advancements reduce the labor required for a task. The labor demand function for workers is given by:
\[ L = 100 - 2w \]
where $L$ is the number of workers employed, and $w$ is the hourly wage.

With automation, the demand function changes to:
\[ L = 80 - 2w \]

\textbf{Questions:}
\begin{enumerate}
    \item If the initial wage is \$20/hour, how many workers are employed before and after automation?
    \item Determine the wage at which no workers are employed before and after automation.
    \item Discuss the implications of this change for workers and the firm.
\end{enumerate}

\section*{Exercise 3: Skill Premium and Automation}
\textbf{Scenario:} A firm adopts a new AI technology that complements high-skilled labor but replaces low-skilled labor. Wages for high-skilled workers rise from \$30/hour to \$40/hour, while wages for low-skilled workers drop from \$20/hour to \$15/hour.

\textbf{Questions:}
\begin{enumerate}
    \item Calculate the percentage change in wages for both high-skilled and low-skilled workers.
    \item If the firm employs 50 high-skilled workers and 100 low-skilled workers before automation, what is the total wage bill before and after automation?
    \item Analyze how the adoption of AI affects income inequality within the firm.
\end{enumerate}

\section*{Exercise 4: Technological Adoption and Output Growth}
\textbf{Scenario:} A company invests \$1 million in new automation technology, which increases its annual output from 50,000 units to 80,000 units. The cost per unit decreases from \$20 to \$15 due to efficiency gains.

\textbf{Questions:}
\begin{enumerate}
    \item What is the percentage increase in output after adopting the technology?
    \item Calculate the change in total revenue if the selling price per unit remains \$25.
    \item Discuss how technological adoption can impact market competitiveness and prices.
\end{enumerate}

\section*{Exercise 5: Job Polarization}
\textbf{Scenario:} In an economy, job distribution shifts due to automation. The initial job shares across sectors are as follows:
\begin{itemize}
    \item Low-skill jobs: 40\%
    \item Middle-skill jobs: 40\%
    \item High-skill jobs: 20\%
\end{itemize}

After automation, the shares change:
\begin{itemize}
    \item Low-skill jobs: 45\%
    \item Middle-skill jobs: 25\%
    \item High-skill jobs: 30\%
\end{itemize}

\textbf{Questions:}
\begin{enumerate}
    \item Calculate the percentage point change in job shares for each skill level.
    \item If the total labor force is 10 million, how many jobs are there in each category before and after automation?
    \item Discuss how automation leads to job polarization and its implications for income inequality.
\end{enumerate}

\section*{Exercise 6: Automation's Effect on GDP Growth}
\textbf{Scenario:} A country’s GDP is modeled as:
\[ GDP = A \times L^{0.6} \times K^{0.4} \]
where $A$ is the technology level, $L$ is labor input, and $K$ is capital input.

Initially:
\begin{itemize}
    \item $A = 1$, $L = 1,000$, $K = 500$.
\end{itemize}

After automation:
\begin{itemize}
    \item $A = 1.2$, $L = 800$, $K = 600$.
\end{itemize}

\section*{Notes on the Exercises}

\subsection*{Exercise 1: Labor Productivity Growth}
\begin{itemize}
    \item \textbf{Labor Productivity:} Automation typically increases productivity by enabling workers to focus on higher-value tasks or substituting manual processes.
    \item \textbf{Implications:} Higher productivity can lead to increased competitiveness and wages for skilled workers but may cause job displacement for others.
\end{itemize}

\subsection*{Exercise 2: The Impact of Automation on Employment}
\begin{itemize}
    \item \textbf{Labor Demand:} Automation shifts the labor demand curve, reducing employment at the same wage level.
    \item \textbf{Implications:} Firms benefit from efficiency gains, but displaced workers may face challenges unless they transition to other roles.
\end{itemize}

\subsection*{Exercise 3: Skill Premium and Automation}
\begin{itemize}
    \item \textbf{Skill Premium:} Automation often complements high-skilled labor and substitutes low-skilled labor, increasing wage inequality.
    \item \textbf{Implications:} The rise in wage disparity can lead to broader societal challenges if unchecked.
\end{itemize}

\subsection*{Exercise 4: Technological Adoption and Output Growth}
\begin{itemize}
    \item \textbf{Output Growth:} Automation increases production and reduces per-unit costs, enhancing competitiveness.
    \item \textbf{Implications:} Consumers benefit from lower prices, and firms capture larger market shares if quality is maintained.
\end{itemize}

\subsection*{Exercise 5: Job Polarization}
\begin{itemize}
    \item \textbf{Job Distribution:} Automation increases demand for high-skill and low-skill jobs but reduces middle-skill job opportunities.
    \item \textbf{Implications:} This polarization may widen income inequality and requires targeted policy interventions.
\end{itemize}

\subsection*{Exercise 6: Automation's Effect on GDP Growth}
\begin{itemize}
    \item \textbf{GDP Growth:} Technological improvements boost GDP by increasing the productivity of labor and capital.
    \item \textbf{Implications:} While GDP rises, reduced labor input may create challenges such as unemployment or skill mismatches.
\end{itemize}

\end{document}
